\subsection{Clifford group.}

In this section, we assume that A is a field, that E has finite dimension m, and that Q is nondegenerate. We identify E with its canonical image in C(Q).

DEFINITION 2. --- We call the Clifford group of Q (resp. the special Clifford group of Q) the multiplicative group of invertible elements s of C(Q) (resp. C \(^{+}\) (Q)) such that sEs \(^{-1}\) = E.

In this section, we denote by G and G⁺ the Clifford group and the special Clifford group of Q. It is clear that G⁺ = G ∩ C⁺(Q).

THEOREM 4. --- Set, for \(s \in G\) and \(x \in E\) , \(\varphi(s). x = sxs^{-1}\) .

\begin{enumerate}
\def\labelenumi{\alph{enumi})}
\item
  The map \(\varphi\) is a homomorphism from G into the orthogonal group \(\mathbf{O}(\mathrm{Q})\) of Q and its kernel is the set of invertible elements of the center Z of C(Q).
\item
  The set \(E \cap G\) is the set of nonsingular vectors of E; for \(x \in E \cap G, -\varphi(x)\) is the symmetry with respect to the hyperplane orthogonal to x.
\item
  If dim(E) is even, we have \(\varphi(G) = \mathbf{O}(Q)\) , \(\varphi(G^{+})\) has index 2 in \(\mathbf{O}(Q)\) if \(E \neq \{0\}\) , and is equal to \(\mathbf{SO}(Q)\) if A has characteristic \(\neq 2\) .
\item
  If dim(E) is odd (which implies that A has characteristic ≠ 2), we have φ(G) = φ(G⁺) = SO(Q).
\end{enumerate}

Indeed, we have \(\mathrm{Q}(sxs^{-1}) = (sxs^{-1})^2 = sx^2s^{-1} = \mathrm{Q}(x)\) for \(s \in G\) and \(x \in E\) , which shows that \(\varphi(s) \in \mathbf{O}(Q)\) . For \(\varphi(s) = 1\) , it is necessary and sufficient that s commute with the elements of E, that is, belong to the center Z of C(Q). This proves a).

For an element \(x\) of E to belong to G, it must be invertible, that is, it must be a non-singular vector (since \(x^2 = Q(x)\) ). If this is so, we have \(x^{-1} = Q(x)^{-1}x\) , whence, for every \(y \in E\) ,

\[
x y x ^ {- 1} = \mathrm{Q} (x) ^ {- 1} x y x = \mathrm{Q} (x) ^ {- 1} x (\Phi (x, y) - x y) = - (y - \Phi (x, y) \mathrm{Q} (x) ^ {- 1} x)
\]

this proves b) (§ 6, no. 4).

Lemma 5. --- Every element s of G is of the form \(zs'\) , where z is an invertible element of Z and \(s'\) belongs to G ∩ C \(^{+}\) (Q) or to G ∩ C \(^{-}\) (Q); the subgroup G \(^{+}\) has index 2 in G when E ≠ \{0\}.

The second assertion follows evidently from the first, since the nonsingular vectors belong to \(\mathrm{G}\cap\mathrm{C}^{-}(\mathrm{Q})\) . Suppose first that dim(E) is even, and let \(s=t'+t''\) , with \(t'\in\mathrm{C}^{+}(\mathrm{Q})\) and \(t''\in\mathrm{C}^{-}(\mathrm{Q})\) ; by definition we have \(sx=(\varphi(s).x)s\) for all \(x\in E\) ; since \(t'x\) and \((\varphi(s).x)t'\) (resp. \(t''x\) and \((\varphi(s).x)t'')\) are odd (resp. even) elements, we have \(t'x=(\varphi(s).x)t'\) , whence \(s^{-1}t'x=xs^{-1}t'\) for all \(x\in E\) . From this we conclude that \(s^{-1}t'\in Z\) and since dim(E) is even, \(Z=A\) (no. 4, cor. of th. 2) , hence \(t'=as\) , where \(a\in A\) . If \(a\neq0\) , we therefore have \(s=a^{-1}t'\) and \(t'\in G\cap C^{+}(Q)\) ; if a=0, \(s=t''\in G\cap C^{-}(Q)\) and the lemma is proved in this case. If dim(E) is odd, A has characteristic \(\neq2\) , so for every \(s\in G\) , \(\varphi(s)\) is a product of symmetries with respect to nonsingular vectors \(x_{i} (i = 1,2,\ldots,h)\) ( \(\S 6, n^{o} 4\) , prop. 5); if we set \(s' = x_{1}x_{2}\ldots x_{h}\) , one has \(\varphi(s) = \varphi(s')\) , hence \(s = zs'\) , where \(z \in \mathbf{Z}\) , and \(s'\) belongs to \(C^{+}(Q)\) or to \(C^{-}(Q)\) according as \(h\) is even or odd.

Suppose dim(E) is even. Since every element \(u\) of \(\mathbf{O}(Q)\) extends in one and only one way to an automorphism \(\overline{u}\) of C(Q) (prop. 1), and since C(Q) is central simple (th. 2), \(\overline{u}\) is an inner automorphism (chap.~VIII, § 10, no. 1, th. 1). There therefore exists an element \(s\) of G such that \(\varphi(s) = u\) On the other hand, the center of C(Q) is contained in C \(^{+}\) , which implies that \(\varphi(G)/\varphi(G^{+})\) is isomorphic to G/G \(^{+}\) , hence that \(\varphi(G^{+})\) has index 2 in \(\varphi(G) = \mathbf{O}(Q)\) if E ≠ \{0\}. This proves the first two assertions of \(c\) ).

Finally, suppose that A has characteristic ≠ 2. Then every element u of O(Q) is a product of symmetries with respect to hyperplanes orthogonal to nonsingular vectors \(x_{i} (i = 1, \ldots, h) (\S 6, n^{o} 4, \text{prop. } 5)\) ; we therefore have \(u = (-1)^{h} \varphi(x_{1} \ldots x_{h})\) and det \((u) = (-1)^{h}\) . For u to belong to SO(Q), it is necessary and sufficient that h be even, which shows that \(\varphi(G^{+}) \supset \mathbf{SO}(Q)\) . Since SO(Q) has index 2 in O(Q) when \(E \neq \{0\}\) , one has \(\varphi(G^{+}) = \mathbf{SO}(Q)\) If E is of even dimension, which completes the proof of c). On the other hand, if the dimension of E is odd, \(\varphi(G)\) does not contain the orthogonal transformation \(x \to -x\) ; indeed, the latter extends to the principal automorphism \(\alpha\) of C(Q) (no. 1), and \(\alpha\) is not an inner automorphism since the center Z of C(Q) contains a nonzero element of C⁻(Q) (th. 3). We therefore have \(\varphi(G) \neq \mathbf{O}(Q)\) , and, since \(\varphi(G) \supset \varphi(G^{+}) \supset \mathbf{SO}(Q)\) and since SO(Q) is of index 2 in O(Q), we have \(\varphi(G) = \varphi(G^{+}) = \mathbf{SO}(Q)\) This proves d). QED.

The subgroup \(\varphi(G^{+})\) of \(\mathbf{O}(Q)\) , which is of index 2 if \(E \neq \{0\}\) , is called the group of rotations of E, and its elements are called rotations; it is denoted by \(\mathbf{O}^{+}(Q)\) Note that, if A is not of characteristic 2, we have \(\mathbf{O}^{+}(Q) = \mathbf{S}\mathbf{O}(Q)\) (cf.~exerc. 9).

PROPOSITION 4. — Let β be the principal antiautomorphism of C(Q) (no. 1). For every \(s \in G^{+}\) , \(\beta(s)s\) is a scalar. The map N : \(s \to \beta(s)s\) is a homomorphism of \(G^{+}\) in the multiplicative group A* of nonzero elements of A.

Indeed, for \(s \in G^{+}\) , one has \(sEs^{-1} = E\) , whence \(\beta(s)^{-1}E\beta(s) = E\) , which shows that \(\beta(s) \in G^{+}\) . Since \(sx = (\varphi(s).x)s\) for all \(x \in E\) , one has \(x\beta(s) = \beta(sx) = \beta(s)(\varphi(s).x)\) and consequently \(\beta(s)sx = \beta(s)(\varphi(s).x)s = x\beta(s)s\) , which implies that \(\beta(s)s\) belongs to the center of C(Q). Since, moreover, \(\beta(s)s\) belongs to C+(Q), \(\beta(s)s\) is a scalar (Ths. 2 and 3). Finally we have \(\beta(st)st = \beta(t)\beta(s)st = \beta(s)s\beta(t)t\) , that is, N(st) = N(s)N(t) for s, t in G+. QED.

The scalar \(N(s) = \beta(s)s (s \in G^{+})\) is called the spinorial norm of s. One denotes by \(G_{0}^{+}\) and one calls the reduced Clifford group the kernel of the homomorphism N. The image \(\varphi(G_{0}^{+})\) is denoted \(\mathbf{O}_{0}^{+}(Q)\) and is called the reduced orthogonal group of Q. Since the kernel of the restriction of \(\varphi\) to \(G^{+}\) is the set of even invertible elements of the center of C(Q) (Th. 4) and hence is identified with A* (Ths. 2 and 3), \(\varphi(G^{+})/\mathbf{O}_{0}^{+}(Q)\) is isomorphic to \(G^{+}/A^{*}G_{0}^{+}\) , hence also with \(N(G^{+})/N(A^{*})\) , and in particular is commutative. It is clear that \(N(A^{*})\) is the subgroup \((A^{*})^{2}\) of squares of elements of A*. If Q has index \textgreater0, there exist, for every \(a \in A^{*}\) , two elements x and y of E such that \(Q(x) = a\) and \(Q(y) = 1 (\S 4, n^{\circ} 2, \text{prop. } 4)\) ; since \(xy \in G^{+}\) and \(N(xy) = Q(x)Q(y) = a\) ; this shows that \(N(G^{+}) = A^{*}\) , hence that \(\varphi(G^{+})/\mathbf{O}_{0}^{+}(Q)\) is isomorphic to \(A^{*}/(A^{*})^{2}\) .

Exercises. --- ¶ 1) Prove Corollaries 3 and 4 of Theorem 1 of No. 3 when E₁ and E₂ are any two complementary submodules of E. (First establish Cor. 4: to do this, show that the tensor product C(Q₁) ⊗ C(Q₂) is endowed with an algebra structure by the sign convention made in the statement of Cor. 4, and that this algebra S solves the same universal mapping problem as C(Q); for this, consider, for every linear map f of E into an algebra D over A such that (f(x))² = Q(x)·1, the homomorphism \(\bar{f}_i\) of C(Qᵢ) into D such that \(\bar{f}_i(1) = 1\) , \(\bar{f}_i(\rho_{Q_i}(x_i)) = f(x_i)\) for \(x_i \in \mathrm{E}_i\) ( \(i = 1, 2\) ) and prove that there exists a homomorphism \(\bar{f}\) of the algebra S into D such that

\[
\bar {f} (z _ {1} \otimes z _ {2}) = \bar {f} _ {1} (z _ {1}) \bar {f} _ {2} (z _ {2})
\]

for \(z_{i}\in\mathrm{C}(\mathrm{Q}_{i})\) ( \(i=1,2\) ). To then prove Cor. 3, consider the quadratic form Q' which is the external direct sum of Q \(_{1}\) and Q \(_{2}\) (§ 3, No. 4), and note that one has Q'(x)=Q(x)+F(x,x), F being the bilinear form defined by F(x \(_{1}\) +x \(_{2}\) , y \(_{1}\) +y \(_{2}\) ) = -Φ(x \(_{1}\) , y \(_{2}\) ) for x \(_{i}\) , y \(_{i}\) in E \(_{i}\) ( \(i=1,2\) ); then use prop. 3.)

¶ 2) Suppose that E is the direct sum of two orthogonal submodules E \(_{1}\) , E \(_{2}\) and denote by Q \(_{i}\) the restriction of Q to E \(_{i}\) (i = 1, 2);

suppose moreover that there exists \(u \in \mathrm{C}^{+}(\mathrm{Q}_{2})\) such that \(u^{2} = a.1\) , \(a\) being invertible in A, and that \(u\rho_{\mathrm{Q}_{2}}(x_{2}) = -\rho_{\mathrm{Q}_{2}}(x_{2})\) for all \(x_{2} \in \mathrm{E}_{2}\) . Show that there exists an isomorphism \(\varphi\) of C(Q) onto the tensor product \(\mathrm{C}(a\mathrm{Q}_{1}) \otimes \mathrm{C}(\mathrm{Q}_{2})\) (as defined in Chap.~III, § 3, no. 1), having the property that for every \(x = x_{1} + x_{2}\) ( \(x_{i} \in \mathrm{E}_{i}, i = 1, 2\) ), one has

\[
\varphi (\rho_ {q} (x)) = \rho_ {a q _ {1}} (x _ {1}) \otimes u ^ {- 1} + 1 \otimes \rho_ {q _ {2}} (x _ {2}).
\]

(Prove the existence of the homomorphism \(\varphi\) as a consequence of the universal property of C(Q). On the other hand, there is a homomorphism \(g_{1}\) of C(aQ \(_{1}\) ) into C(Q) such that \(g_{1}(\rho_{aQ_{1}}(x_{1})) = h_{2}(u)\rho_{Q_{1}}(x_{1})\) , denoting by \(h_{2}\) the canonical homomorphism of C(Q \(_{2}\) ) into C(Q); then observe that \(g_{1}(z_{1})\) and \(h_{2}(z_{2})\) commute for \(z_{1} \in\) C(aQ \(_{1}\) ) and \(z_{2} \in\) C(Q \(_{2}\) ), and deduce from this the existence of an inverse homomorphism of \(\varphi\) .)

In which case does this result apply when A is a field of characteristic ≠ 2 (use Remark 2 following the cor. of th. 2)?

\begin{enumerate}
\def\labelenumi{\arabic{enumi})}
\setcounter{enumi}{2}
\tightlist
\item
  \begin{enumerate}
  \def\labelenumii{\alph{enumii})}
  \tightlist
  \item
    With the notations of no. 2, let \(x_{i}\) ( \(1 \leqslant i \leqslant n\) ) be elements of E such that \(f(x_{i}) = 0\) ; show that one has \(i_{j}(x_{1} \otimes x_{2} \otimes \cdots \otimes x_{n}) = 0\) . In particular, if F is a bilinear form on E such that F( \(x_{i}, x_{j}\) ) = 0 for \(i > j\) , one has \(i_{x_{1}}^{\mathrm{F}}(x_{2} \otimes x_{3} \otimes \cdots \otimes x_{n}) = 0\) .
  \end{enumerate}
\end{enumerate}

\begin{enumerate}
\def\labelenumi{\alph{enumi})}
\setcounter{enumi}{1}
\tightlist
\item
  With the notations of Prop. 3 of no. 3, let \(x_i\) ( \(1 \leqslant i \leqslant n\) ) be elements of E such that \(F(x_i, x_j) = 0\) for \(i > j\) . Show that we have \(\overline{\lambda}_{\mathrm{F}}(\rho_{\mathrm{Q}'}(x_1) \ldots \rho_{\mathrm{Q}'}(x_n)) = \rho_{\mathrm{Q}}(x_1) \ldots \rho_{\mathrm{Q}}(x_n)\) (using \(a\) ) and formula (10) of no. 2).
\item
  We assume that A is a field of characteristic \(\neq 2\) ; for every quadratic form Q on E, let \(\mu_{\mathrm{Q}}\) be the map \(\overline{\lambda}_{\mathrm{F}}\) of C(Q) onto \(\wedge\) E corresponding to \(F(x, y) = \frac{1}{2}\Phi(x, y)\) . Show that if the vectors \(x_i\) ( \(1 \leqslant i \leqslant n\) ) are pairwise orthogonal, we have
\end{enumerate}

\[
\mu_ {\mathrm{Q}} (x _ {1} x _ {2} \dots x _ {n}) = x _ {1} \wedge x _ {2} \wedge \dots \wedge x _ {n}.
\]

Deduce that if \(y_{i}\) ( \(1 \leqslant i \leqslant n\) ) are \(n\) arbitrary vectors, we have

\[
n! \mu_ {0} ^ {- 1} (y _ {1} \wedge y _ {2} \wedge \dots \wedge y _ {n}) = \sum_ {\sigma \in \mathfrak {S} _ {n}} \varepsilon_ {\sigma} y _ {\sigma (1)} y _ {\sigma (2)} \dots y _ {\sigma (n)}.
\]

\begin{enumerate}
\def\labelenumi{\arabic{enumi})}
\setcounter{enumi}{3}
\tightlist
\item
  Let \(C_{h}\) be the submodule of C(Q) generated by the products of k elements of \(\rho_{Q}(E)\) for \(0 \leqslant k \leqslant h\) . Show that the map
\end{enumerate}

\[
(x _ {1}, \dots , x _ {h}) \rightarrow \rho_ {\mathrm{Q}} (x _ {1}) \dots \rho_ {\mathrm{Q}} (x _ {h})
\]

defines by passing to quotients an alternating multilinear map from \(E^{h}\) into \(C_{h}/C_{h-1}\) . From this we deduce a linear map \(\pi_{h}\) of \(\bigwedge E\) into \(C_{h}/C_{h-1}\) ; show that \(\pi_{h}\) is an isomorphism when E is a free module. If f is an orthogonal transformation, we have \(C(f)\circ\pi_{h}=\pi_{h}\circ(\bigwedge f)\) . When A is a field of characteristic \(\neq2\) , show that \(\pi_{h}\) is the restriction of the map \(\mu_{Q}^{-1}\) defined in Exerc. 3 c).

\begin{enumerate}
\def\labelenumi{\arabic{enumi})}
\setcounter{enumi}{4}
\tightlist
\item
  With the notation of no. 2, consider the map \(i_{f}\) of \(\wedge\) from E into itself. Show that one has
\end{enumerate}

\[
i _ {f} \left(x _ {1} \wedge \dots \wedge x _ {h}\right) = \sum_ {i = 1} ^ {h} (- 1) ^ {i - 1} f \left(x _ {i}\right) \left(x _ {1} \wedge \dots \wedge x _ {i - 1} \wedge x _ {i + 1} \wedge \dots \wedge x _ {h}\right)
\]

and deduce that \(i_{f}\) is none other than the right interior product \(z \to z \sqcup f\) (chap.~III, § 8, no. 4, def. 2).

\begin{enumerate}
\def\labelenumi{\arabic{enumi})}
\setcounter{enumi}{5}
\item
  Show that, if E admits an orthogonal basis \((e_{1}, e_{2})\) for Q, and if we set \(\alpha_{i} = \mathrm{Q}(e_{i})\) For (i = 1, 2), the algebra C(Q) is isomorphic to the algebra of quaternions over A corresponding to the pair \((\alpha_{1}, \alpha_{2})\) .
\item
  Let A be a maximal ordered field, E a vector space of even dimension 2r over A, Q a nondegenerate quadratic form on E, positive or negative. Show that, if Q is positive, C(Q) is isomorphic to a matrix algebra over A when \(r(r-1)/2\) is even, to a matrix algebra over the quaternion algebra over A when \(r(r-1)/2\) is odd. If Q is negative, C(Q) is isomorphic to a matrix algebra over A when \(r(r+1)/2\) is even, to a matrix algebra over the quaternion algebra over A when \(r(r+1)/2\) is odd (use Exercises 2 and 6). What is the structure of C⁺(Q) in these different cases?
\item
  Let A be the ring \(\mathbf{Z}/(4)\) a left A-module \(\mathbf{Z}/(2)\) ; if we set \(Q(0)=0\) , \(Q(u)=1\) for the unique element \(u\neq0\) of E, Q is a quadratic form on E. Show that the A-modules C(Q) and \(\wedge\) They are not isomorphic (one will prove that C(Q) is isomorphic to the direct sum of two modules isomorphic to E).
\end{enumerate}

¶ 9) Suppose A is a field of characteristic 2, E a vector space of finite even dimension 2r, and Q a nondegenerate quadratic form on E.

\begin{enumerate}
\def\labelenumi{\alph{enumi})}
\tightlist
\item
  Let \((e_i)\) a symplectic basis (§ 5, no. 1) of E for the alternating bilinear form \(\Phi\) associated with Q. Show that the element
\end{enumerate}

\[
z = e _ {1} e _ {2} + e _ {3} e _ {4} + \dots + e _ {2 r - 1} e _ {2 r}
\]

of C \(^{+}\) (Q), together with the identity element, forms a basis of the center Z of C \(^{+}\) For Z to be a direct sum of two fields, it is necessary and sufficient that the element

\[
\Delta (Q) = Q \left(e _ {1}\right) Q \left(e _ {2}\right) + Q \left(e _ {3}\right) Q \left(e _ {4}\right) + \dots + Q \left(e _ {2 r - 1}\right) Q \left(e _ {2 r}\right)
\]

(called the pseudo-discriminant of Q with respect to the symplectic basis \((e_{i})\) ) is of the form \(\lambda^{2} + \lambda\) , with \(\lambda \in A\) .

\begin{enumerate}
\def\labelenumi{\alph{enumi})}
\setcounter{enumi}{1}
\tightlist
\item
  Let \(u\) a similarity for \(\Phi\) ( \(\S 6\) , no. 5) with multiplier \(\mu(u)\) , and let \(Q_{1}(x) = Q(u(x))\) . We set
\end{enumerate}

\[
u (e _ {2 i - 1}) = \sum_ {j = 1} ^ {r} a _ {i j} e _ {2 j - 1} + \sum_ {j = 1} ^ {r} b _ {i j} e _ {2 j},
\]

\[
u (e _ {2 i}) = \sum_ {j = 1} ^ {r} c _ {i j} e _ {2 j - 1} + \sum_ {j = 1} ^ {r} d _ {i j} e _ {2 j},
\]

and \(\mathrm{Q}(e_{2i - 1}) = \alpha_i, \mathrm{Q}(e_{2i}) = \beta_i (1 \leqslant i \leqslant r)\) . Show that we have

\[
\Delta (\mathrm{Q} _ {1}) = (\mu (u)) ^ {2} \Delta (\mathrm{Q}) + (\mathrm{D} (u)) ^ {2} + \mu (u) \mathrm{D} (u)
\]

where

\[
\mathbf {D} (u) = \sum_ {i, j} (\alpha_ {j} a _ {i j} c _ {i j} + \beta_ {j} b _ {i j} d _ {i j} + b _ {i j} c _ {i j})
\]

(Dickson invariant of u with respect to the basis \((e_{i})\) ). (Note that the element

\[
(\mu (u)) ^ {- 1} (u (e _ {1}) u (e _ {2}) + u (e _ {3}) u (e _ {4}) + \dots + u (e _ {2 r - 1}) u (e _ {2 r}))
\]

belongs to Z.) For \(u\) to be a similarity for Q (§ 4, exerc. 9) with multiplier \(\mu(u)\) , it is necessary and sufficient that D(u) = 0 or D(u) = \(\mu(u)\) ; the similarities for Q such that D(u) = 0 are called direct.

\begin{enumerate}
\def\labelenumi{\alph{enumi})}
\setcounter{enumi}{2}
\tightlist
\item
  Show that, if v is a similarity for \(\Phi\) , u a similarity for Q, one has
\end{enumerate}

\[
\mathrm{D} (u v) = \mu (v) \mathrm{D} (u) + \mu (u) \mathrm{D} (v)
\]

(consider the Dickson invariant of \(\varrho\) with respect to the symplectic basis formed by the \(u(e_{2i-1})\) and \((\mu(u))^{-1}u(e_{2i})\) for \(1 \leqslant i \leqslant r\) ). Deduce from this that the direct similarities for Q form a distinguished subgroup of index 2 in the group of similarities for Q.

\begin{enumerate}
\def\labelenumi{\alph{enumi})}
\setcounter{enumi}{3}
\item
  If \(u\) is the symmetry with respect to the hyperplane orthogonal to a nonsingular vector in E, show that \(\mathrm{D}(u)=1\) . Deduce from this that the group \(\varphi(\mathrm{G}^{+})\) (notations of no. 5) is the group \(\mathbf{SO}(\mathrm{Q})\) defined in exerc. 28 c) of § 6.
\item
  Let \(u \in \mathbf{O}(Q)\) , and suppose \(\Lambda\) algebraically closed. Show that, for \(u \in \mathbf{SO}(Q)\) , it is necessary and sufficient that the number of elementary divisors of the module \(E_u\) be even. (With the notations of exerc. 15 of § 5, first note that if \(p, q\) are two distinct irreducible factors of the minimal polynomial of \(u\) , the number of indecomposable submodules of which \(G(p, q)\) is the direct sum is equal to the number of indecomposable submodules of which \(G(q, p)\) is the direct sum. Note on the other hand that \(G(p, p) = \{0\}\) except for \(p = X - 1\) . Finally prove that one may restrict to the case where \(E_u\) is equal to \(G(p, p)\) (with \(p = X - 1\) ) and is indecomposable. There is then a symplectic basis ( \(e_i\) ) of E such that \(e_1, e_3, \ldots, e_{2k-1}\) form a basis of \((p(u))^{2r-k}(E)\) for \(1 \leqslant k \leqslant r\) ; show that \(e_1, e_3, \ldots, e_{2r-3}\) are singular vectors, and conclude that \(D(u) = 1\) .)
\end{enumerate}

¶ 10) Suppose A is a field, E a finite-dimensional vector space, Q a degenerate quadratic form on E; let M be a subspace complementary to E \(^{0}\) in E, let B be the (semisimple) Clifford algebra of the restriction of Q to M.

\begin{enumerate}
\def\labelenumi{\alph{enumi})}
\item
  We first assume that A has characteristic ≠ 2. Let L be the Clifford algebra of the restriction of Q to E⁰ (isomorphic to ∧ E⁰), and R₀ its radical (the ideal generated in L by E⁰, of codimension 1 in L); show that the radical R of C(Q) is obtained (up to isomorphism) by defining the algebra structure on \(B \otimes_{A} R_{0}\) as in cor. 4 of th. 1, that C(Q)/R is isomorphic to B, and that C(Q) is the direct sum of B and R.
\item
  Assume A has characteristic 2. Let F be the subspace of \(\mathbf{E}^0\) consisting of singular vectors \(x \in \mathbf{E}^0\) and let N be a complement of F with respect to \(\mathbf{E}^0\) . If \((a_i)_{1 \leqslant i \leqslant d}\) is a basis of N, and \(\mathrm{Q}(a_i) = \alpha_i\) the elements \(\alpha_i^{1/2}\) in an algebraic closure of A, are linearly independent over A. Let \((\alpha_i^{1/2})_{1 \leqslant i \leqslant e}\) be a 2-basis of the field \(\mathrm{A}_1 = \mathrm{A}(\alpha_1^{1/2}, \ldots, \alpha_d^{1/2})\) on A (Chap.~V, § 8, exerc. 1), and let \(h = \dim F\) . If \(\mathrm{B}_1\) is the central simple algebra \(\mathrm{B} \otimes_{A} \mathrm{A}_1\) , show that \(\mathrm{C(Q)}\) is isomorphic to the algebra \(\mathrm{B}_1 \otimes_{A_1} \mathrm{L}_1\) , where \(\mathrm{L}_1\) is the exterior algebra of a vector space of dimension \(h + d - e\) on \(\mathrm{A}_1\) . If \(\Re_1\) is the radical of \(\mathrm{L}_1\) (of codimension 1 (over \(\mathrm{A}_1\) ) in \(\mathrm{L}_1\) ), the radical \(\Re\) of \(\mathrm{C(Q)}\) is isomorphic to the algebra \(\mathrm{B}_1 \otimes_{A_1} \Re_1\) , \(\mathrm{C(Q)/R}\) is isomorphic to \(\mathrm{B}_1\) , and \(\mathrm{C(Q)}\) is a direct sum of \(\mathrm{B}_1\) and of \(\Re\) .
\item
  The hypotheses being those of \(b\) ; one assumes moreover that \(h=0\) , \(d=1\) , \(e=0\) (which implies that dim M is even). If Q \(_{0}\) is the restriction of Q to M, show that C \(^{+}\) (Q) is isomorphic to C(Q \(_{0}\) ).
\end{enumerate}

11) a) Under the hypotheses and notation of no. 5, show that N(G⁺) is the subgroup of A* generated by the products Q(x)Q(y), where x and y range over the set of non-isotropic vectors of E. (Reduce to the case A ≠ F₂; use prop. 5 and exerc. 28 c) of § 6, as well as exerc. 9 d) of § 9.) Case where Q has index ≥ 1.

\begin{enumerate}
\def\labelenumi{\alph{enumi})}
\setcounter{enumi}{1}
\item
  We assume in addition that \(\mathbf{A} \neq \mathbf{F}_2\) that E has dimension \(n \geqslant 2\) and that Q has index \(>0\) . Show that \(\mathbf{O}_0^+(\mathbf{Q})\) is the commutator subgroup of the group \(\mathbf{O}(\mathbf{Q})\) . (Reduce to the case \(n = 2\) using Exercises 17 c) and 28 e) of § 6.)
\item
  We keep the hypotheses of \(b\) , and we suppose in addition that A has characteristic \(\neq 2\) and that E is of even dimension. For the automorphism \(x \to -x\) of E belong to \(\mathbf{O}_0^+(\mathbf{Q})\) it is necessary and sufficient that the discriminant of Q (with respect to any basis of E) be a square.
\end{enumerate}

Let a be an invertible element of A; show that there exists one and only one isomorphism \(\theta_{a}\) of C \(^{+}\) (Q) onto C \(^{+}\) (aQ) such that

\[
\theta_ {a} (\rho (x) \rho (y)) = a ^ {- 1} \rho_ {1} (x) \rho_ {1} (y)
\]

for all \(x, y\) in E ( \(\rho\) and \(\rho_{1}\) denoting the canonical maps from E into C(Q) and C(aQ), respectively). Deduce that if A is a field of characteristic \(\neq 2\) , E a vector space of odd dimension and Q a non-degenerate quadratic form, C(Q) and C(aQ) are isomorphic (cf.~exerc. 7).

\begin{enumerate}
\def\labelenumi{\alph{enumi})}
\setcounter{enumi}{1}
\tightlist
\item
  Assume that A is a field, E a vector space of even dimension 2r \textgreater{} 0, Q a nondegenerate quadratic form. Let u be a similarity relative to Q; show that there exists one and only one A-automorphism \(\overline{u}\) of the algebra \(\mathrm{C}^{+}(\mathrm{Q})\) such that for \(1 \leqslant h \leqslant r\) and \(x_{i} \in E (1 \leqslant i \leqslant 2h)\) one has
\end{enumerate}

\[
\overline {{{u}}} (x _ {1} x _ {2} \dots x _ {2 h}) = \mu^ {- h} u (x _ {1}) u (x _ {2}) \dots u (x _ {2 h})
\]

where \(\mu\) denotes the multiplier of \(u\) . For \(\bar{u}\) to be an inner automorphism, it is necessary and sufficient that \(u\) be a direct similarity (§ 6, no. 5 and § 9, exerc. 9 b)). One assumes in addition \(r \geqslant 2\) ; then, for \(\bar{u}\) to be the identity, it is necessary and sufficient that \(u\) be a homothety.

Show that the automorphism \(\bar{u}\) of C \(^{+}\) (Q) is the restriction to C \(^{+}\) (Q) of an inner automorphism of C(Q).

¶ 13) Suppose that A is a field of characteristic ≠ 2, E a vector space of even dimension 2r. \textgreater{} 0, Q a nondegenerate quadratic form. We denote by \(E_{h}\) the inverse image of \(\Lambda E\) under the isomorphism \(\mu_{Q}\) of exerc. 3 c), so that \(E_{h}\) is the subspace of C(Q) generated by the products \(x_{1}x_{2}\ldots x_{h}\) , where the \(x_{i}\) ( \(1 \leqslant i \leqslant h\) ) are pairwise orthogonal.

\begin{enumerate}
\def\labelenumi{\alph{enumi})}
\tightlist
\item
  Let \(\alpha\) an element of A, \(s\) an element of E \(_2\) . Show that, if \((\alpha + s)^2 \in \mathbf{A}\) , then either \(s = 0\) , or else \(\alpha = 0\) and \(s = xy\) , where \(x\) and \(y\) are two orthogonal vectors. (Note that if \(t = \mu_Q(\alpha + s)\) , \(t \wedge t\)
\end{enumerate}

necessarily belongs to A + ∧ E, by expressing s using an orthogonal basis of E; deduce that one necessarily has t = β + x ∧ y with β ∈ A, x, y in E, using § 5, no. 1, cor. 2 of th. 1).

\begin{enumerate}
\def\labelenumi{\alph{enumi})}
\setcounter{enumi}{1}
\item
  Let \((x, y)\) , \((u, v)\) two pairs of orthogonal non-isotropic vectors in E, \(\mathrm{P}_{xy}\) and \(\mathrm{P}_{uv}\) the planes \(Ax + Ay\) , \(Au + Av\) respectively. For one to have \((xy)(uv) = -(uv)(xy)\) in C(Q), it is necessary and sufficient that \(\mathrm{P}_{xy} + \mathrm{P}_{uv}\) be a non-isotropic subspace of dimension 3, in which \(\mathrm{P}_{xy}\) and \(\mathrm{P}_{uv}\) are weakly orthogonal (§ 3, exerc. 11).
\item
  Let g be an A-automorphism of \(\mathrm{C}^{+}(\mathrm{Q})\) sending \(E_{2}\) into itself; show that there exists a similarity u relative to Q such that \(g = \bar{u}\) (exerc. 12 b)). (Let \((e_{i})_{1 \leqslant i \leqslant 2r}\) be an orthogonal basis of E; using a) and b), show that there exists an orthogonal basis \((x_{i})\) of E such that \(g(e_{1}e_{i}) = x_{1}x_{i}\) for \(2 \leqslant i \leqslant 2r\) .)
\end{enumerate}

\begin{enumerate}
\def\labelenumi{\arabic{enumi})}
\setcounter{enumi}{13}
\tightlist
\item
  We assume that A is a field of characteristic \(\neq 2\) E a vector space of dimension \(n\) on A, Q and Q₁ two non-degenerate quadratic forms on E; let \(\tilde{\Delta}\) (resp. \(\tilde{\Delta}_{1}\) ) be the class of the discriminant \(\Delta\) of Q (resp. of the discriminant \(\Delta_{1}\) of Q \(_{1}\) ) with respect to a basis of E, in the group A \(^{*}\) /(A \(^{*}\) )², a class that does not depend on the basis chosen in E. We assume \(n = 2\) .
\end{enumerate}

\begin{enumerate}
\def\labelenumi{\alph{enumi})}
\item
  Show that if \(\tilde{\Delta} = \tilde{\Delta}_{1}\) and if the Clifford algebras C(Q) and C(Q \(_{1}\) ) are isomorphic, then Q and Q \(_{1}\) are equivalent (consider on C(Q) the quadratic form \(z \to z\bar{z}\) , where \(z \to \bar{z}\) is the unique involutive antiautomorphism of C(Q) whose set of invariants is the center of C(Q) (cf.~exerc. 6 and chap.~VIII, § 11, exerc. 5 e)); apply Witt's theorem).
\item
  For C(Q) to be isomorphic to the matrix algebra \(\mathbf{M}_{2}(\mathrm{A})\) it is necessary and sufficient that, for at least one \(x \neq 0\) in E, there exists \(y\) , \(z\) in E such that Q(x) + Q(y)Q(z) = 0; if this is the case, this property holds for every \(x \neq 0\) in E.
\end{enumerate}

¶ 15) We keep the notations and general hypotheses of exerc. 14, but we assume n = 3.

\begin{enumerate}
\def\labelenumi{\alph{enumi})}
\item
  Show that \(\mathrm{C}^{+}(\mathrm{Q})\) is isomorphic to a quaternion algebra over A and that \(\beta\) is the antiautomorphism \(z\to \bar{z}\) of this algebra whose set of invariants is the center of \(\mathrm{C}^{+}(\mathrm{Q})\) ; if P is the subspace of \(\mathrm{C}^{+}(\mathrm{Q})\) consisting of pure quaternions (that is, those such that \(z = -\bar{z}\) ; chap.~VIII, § 11, exerc. 6), the restriction to P of the quadratic form \(z\to z\bar{z}\) is equivalent to \(\lambda Q\) , where \(\lambda \in A\) . Deduce that, in order for \(C^{+}(Q)\) to be a field, it is necessary and sufficient that \(Q\) be of index 0.
\item
  Show that if \(\tilde{\Delta} = \tilde{\Delta}_{1}\) and if the Clifford algebras C(Q) and C(Q \(_{1}\) ) are isomorphic, then Q and Q \(_{1}\) are equivalent. (First consider the case where - \(\Delta\) is a square in K, and show that in this case C \(^{+}\) (Q) and C \(^{+}\) (Q \(_{1}\) ) are isomorphic; then reason as in exerc. 14 a), using a) and exerc. 6 of chap.~VIII, § 11. In the general case, use exerc. 12 a)).
\item
  Show that the special Clifford group G \(^{+}\) (for the form Q) is identical to the group of invertible elements of C \(^{+}\) (Q). (If ( \(e_{1}\) , \(e_{2}\) , \(e_{3}\) ) is an orthogonal basis of E, and if \(j = e_{1}e_{2}e_{3}\) in C(Q), note that \(x \to xj\) is an isomorphism of vector spaces from E onto P).
\item
  Deduce from a) and c) that if Q has index 1, the group of rotations \(\mathbf{O}^{+}(\mathrm{Q})\) is isomorphic to the projective group \(\mathbf{PGL}_{2}(\mathrm{A})\) (chap.~II, \(2^{\mathrm{e}}\) éd., App. III, no. 6).
\end{enumerate}

We keep the general hypotheses and notations of exerc. 14, but we assume n = 4.

\begin{enumerate}
\def\labelenumi{\alph{enumi})}
\item
  Give an example where \(\tilde{\Delta} = \tilde{\Delta}_{1}\) and where \(C(Q)\) and \(C(Q_{1})\) are isomorphic, but where \(Q\) and \(Q_{1}\) are not equivalent (cf.~exerc. 7).
\item
  Let \((e_i)_{1\leqslant i\leqslant 4}\) an orthogonal basis of E for Q, \(Q_0\) the restriction of Q to the hyperplane H = Ae \(_1\) + Ae \(_2\) + Ae \(_3\) Show that, if Z is the center of C \(^+\) (Q), the algebra C \(^+\) (Q) is isomorphic to the tensor product Z ⊗ \(_A\) C \(^+\) (Q \(_0\) ). For every z ∈ C \(^+\) (Q), one has β(z)z ∈ Z ; for z to belong to the special Clifford group G \(^+\) it is necessary and sufficient that z be invertible and that β(z)z ∈ A.1. Deduce from this that the group O \(_0^+\) (Q) is isomorphic to the quotient by \{1,-1\} of the group of elements z ∈ Z ⊗ \(_A\) C \(^+\) (Q \(_0\) ) such that β(z)z = 1.
\item
  We assume that \(\Delta\) is not a square in A (which implies, by Witt's theorem, that the index of Q is 0 or 1). If \(Q_0'\) is the quadratic form obtained from \(Q_0\) by extension to \(A' = A(\sqrt{\Delta})\) of the field of scalars, deduce from \(b\) ) that \(\mathbf{O}_0^+(Q)\) is isomorphic to \(\mathbf{O}_0^+(Q_0')\) . In particular, if Q has index 1, \(\mathbf{O}_0^+(Q)\) is the commutator subgroup of \(\mathbf{O}(Q)\) and is isomorphic to \(\mathbf{PSL}_2(A')\) (cf.~exerc. 15 \(d\) ) and chap.~III, § 7, exerc. 8).
\item
  We assume that \(\Delta\) is a square in A (which implies, by Witt's theorem, that Q has index 0 or 2) and that \(Q(e_4) = 1\) . If one sets \(j = e_1e_2e_3\)  \(x \in E\) can be written in only one way \(x = \alpha e_4 + jz\) , where \(\alpha \in A\) and \(z\) is a pure quaternion (exerc. 15 a)) in \(L = C^+(Q_0)\) ; if we set \(\psi(x) = \alpha + z\) , \(\psi\) is a vector space isomorphism from E onto L. Let \(Z = Ac' + Ac''\) , where \(c'\) and \(c''\) are the two orthogonal idempotents in Z; every invertible element \(s \in C^+(Q)\) is written in a unique way \(s = uc' + vc''\) , where \(u\) and \(\nu\) belong to L; so that \(s\) belongs to the special Clifford group \(G^+\) , it is necessary and sufficient that \(u\bar{u} = v\bar{v}\) , and we then have \(\psi(sxs^{-1}) = u\psi(x)v^{-1}\) for all \(x \in E\) Deduce from this that the quotient of \(\mathbf{O}_0^+(Q)\) by its center (which is a group with 2 elements, cf.~exerc. 11 c)) is isomorphic to the product \(\mathbf{O}_0^+(Q_0) \times \mathbf{O}_0^+(Q_0)\) ; in particular, if Q has index 2, this quotient group is isomorphic to \(\mathbf{PSL}_2(A) \times \mathbf{PSL}_2(A)\) .
\end{enumerate}

\begin{enumerate}
\def\labelenumi{\arabic{enumi})}
\setcounter{enumi}{16}
\tightlist
\item
  Let K be a commutative field of characteristic ≠ 2, A the field \(\mathrm{K}(\mathrm{X}_{n})_{n\in\mathbb{N}}\) of rational fractions over K, with respect to a countable family of indeterminates (chap.~IV, § 3, n° 1). Let E be a vector space over A, having a countable basis \((e_{n})_{n\in\mathbb{N}}\) , and let Φ be a symmetric bilinear form on E, for which \((e_{n})\) is an orthogonal basis and such that \(\Phi(e_{n}, e_{n}) = \mathrm{X}_{n}\) for all \(n \in N\) . If one sets \(\mathrm{Q}(x) = \Phi(x, x)\) , show that the Clifford algebra C(Q) is a field (cf.~chap.~VIII, § 12, exerc. 14).
\end{enumerate}

\section{Angles}

Throughout this paragraph, A denotes a commutative field of characteristic \(\neq 2\) . Let E be a vector space of dimension 2 over A, and let \(\Phi\) be a nondegenerate symmetric bilinear form on E.

\subsection{Direct similarities in a plane.}

Recall (§ 6, no. 5) that a direct similarity of E is an automorphism u of the vector space E such that \(\Phi(u(x), u(y)) = (\det u)\Phi(x, y)\) for all \(x, y\) in E.

PROPOSITION 1. --- Let A(Φ) be the subalgebra of \(\mathfrak{L}_{\mathbf{A}}(\mathbf{E})\) generated by the direct similarities of E.

\begin{enumerate}
\def\labelenumi{\alph{enumi})}
\item
  Direct similarities are the invertible elements of A(Φ). The algebra A(Φ) is a commutative algebra of degree 2 over A. When E contains no nonzero isotropic vector, A(Φ) is a field, a quadratic extension of A; otherwise it is the direct product of two fields isomorphic to A.
\item
  Let \((e_{1}, e_{2})\) an orthogonal basis of E; set \(\alpha_{i} = \Phi(e_{i}, e_{i})\) (i = 1, 2) and \(\delta = -\alpha_{2}/\alpha_{1}\) Then the matrices of the elements of A( \(\Phi\) ) with respect to this basis are the matrices of the form \(\begin{pmatrix} a & \delta b \\ b & a \end{pmatrix}\) where \(a \in A\) , \(b \in A\) .
\item
  The space E is a free cyclic A(Φ)-module, generated by any non-isotropic vector.
\end{enumerate}

Indeed, let us introduce an auxiliary alternating bilinear form B ≠ 0 on E; then B is non-degenerate. There therefore exists an endomorphism w of E such that \(\Phi(x, y) = \mathrm{B}(w(x), y)\) for all x, y in E. For every invertible endomorphism u of E, we have

\[
\begin{array}{r l} \Phi (u (x), u (y)) & = \mathrm{B} (w u (x), u (y)) = \mathrm{B} (u u ^ {- 1} w u (x), u (y)) \\ & = (\det u) \mathrm{B} (u ^ {- 1} w u (x), y); \end{array}
\]

For u to be a direct similarity, it is necessary and sufficient that one have

\[
(\det u) \mathrm{B} (u ^ {- 1} w u (x), y) = (\det u) \Phi (x, y) = (\det u) \mathrm{B} (w (x), y)
\]

for all \(x, y\) in E; since det \(u \neq 0\) and since B is nondegenerate, this is equivalent to \(u^{-1}wu = w\) , or equivalently to \(uw = wu\) . Take for B the alternating bilinear form whose matrix \(S\) with respect to \((e_1, e_2)\) is \(\begin{pmatrix} 0 & 1 \\ -1 & 0 \end{pmatrix}\) by denoting \(R\) and \(W\) the matrices of \(\Phi\) and of \(w\) with respect to this basis, the relation \(\Phi(x, y) = \mathrm{B}(w(x), y)\) is written, by virtue of formula (47) of § 1, no. 10, \(R = ^t W.S\) ; making this explicit shows that we have \(W = \begin{pmatrix} 0 & \alpha_2 \\ -\alpha_1 & 0 \end{pmatrix}\) . If \(\begin{pmatrix} a & c \\ b & d \end{pmatrix}\) denotes the matrix of \(u\) with respect to \((e_1, e_2)\) , the relation \(uw = wu\) is therefore equivalent to the relations \(b\alpha_2 = -c\alpha_1\) , \(a\alpha_2 = d\alpha_2\) , \(a\alpha_1 = d\alpha_1\) that is, to \(a = d\) and \(c = \delta b\) ; this proves that the matrices of direct similarities are the invertible matrices of the form \(\begin{pmatrix} a & \delta b \\ b & a \end{pmatrix}(a, b\) in A).

Now the endomorphisms of E whose matrices are of the form \(\begin{pmatrix} a & \delta b \\ b & a \end{pmatrix}\) (a, b in A) form a vector subspace of dimension 2 of \(\mathfrak{L}_{\mathrm{A}}(\mathrm{E})\) generated by 1 and by the endomorphism w; since \(w^{2}\) is the homothety with ratio \(-\alpha_{1}\alpha_{2}\) , this subspace is the subalgebra \(\mathrm{A}(\Phi)\) of \(\mathfrak{L}_{\mathrm{A}}(\mathrm{E})\) generated by the direct similarities. The direct similarities are the invertible elements of \(\mathrm{A}(\Phi)\) , that is, those whose matrices satisfy \(a^{2}-\delta b^{2}\neq0\) . The fact that the algebra \(\mathrm{A}(\Phi)\) is commutative is evident. Let us apply to it the results of Chap.~II, §7, no.~7: if δ is not a square in A, that is, if no nonzero vector of E is isotropic, \(\mathrm{A}(\Phi)\) is a field; if, on the contrary, δ is a square in A, that is, if E contains nonzero isotropic vectors, \(\mathrm{A}(\Phi)\) is a direct sum of two fields isomorphic to A. This proves a) and b).

Finally, any non-isotropic vector of E can be taken as

first vector \(e_1\) of an orthogonal basis \((e_1, e_2)\) ( \(\S 6, n^0 1\) ); therefore its transforms \(u(e_1)\) by the elements \(u\) of A( \(\Phi\) ) are the vectors of the form \(ae_1 + be_2\) ( \(a, b\) in A), that is, all vectors of E, since all matrices \(\begin{pmatrix} a & \delta b \\ b & a \end{pmatrix}\) are matrices of elements of A( \(\Phi\) ). In other words, E is a cyclic A( \(\Phi\) )-module, generated by any non-isotropic vector. Moreover, we see that it is a A( \(\Phi\) -module monogenic free, since \(u(e_1) = ae_1 + be_2 = 0\) implies \(a = b = 0\) , hence \(u = 0\) . This proves \(c\) ).

Remarks. --- 1) Let \(\vartheta\) be the similarity with matrix \(\begin{pmatrix}0 & \delta \\ 1 & 0\end{pmatrix}\) with respect to \((e_1, e_2)\) ; the similarity \(w\) introduced in the proof of Prop. 1 is equal to \(-\alpha_1 v\) ; we have \(v^2 = \delta\) . The multiplier of the direct similarity \(u = a + b v\) ( \(a, b\) in A) is equal to the determinant of its matrix \(\begin{pmatrix}a & \delta b \\ b & a\end{pmatrix}\) that is, to \(a^2 - \delta b^2 = (a + b v)(a - b v) = u.\bar{u}\) , denoting by \(\bar{u}\) the conjugate of \(u\) in the algebra A(Φ) (Chap.~II, § 7, No. 7); in other words, the multiplier of \(u\) is the norm N( \(u\) ) of \(u\) in the algebra A(Φ) (ibid.). In particular, for a direct similarity \(u\) to be a rotation, it is necessary and sufficient that N( \(u\) ) = 1; for \(u\) to be a homothety, it is necessary and sufficient that \(u \in A^*\) .

\begin{enumerate}
\def\labelenumi{\arabic{enumi})}
\setcounter{enumi}{1}
\item
  The direct similarities \(u = a + bv\) ( \(a, b\) in A, \(b \neq 0\) ) whose square is a homothety are the homotheties and the scalar multiples \(bv\) of \(v\) , since \((a + bv)^{2} = (a^{2} + \delta b^{2}) + 2abv\) . These latter are none other than the automorphisms of the vector space E which transform every vector into an orthogonal vector; indeed the matrix of such an automorphism is necessarily of the form \(\begin{pmatrix} 0 & c \\ d & 0 \end{pmatrix} (c, d\) in A), and the condition of transforming the vector \(\lambda e_{1} + \mu e_{2}\) as an orthogonal vector is then written \(\lambda \mu(d\alpha_{2} + c\alpha_{1}) = 0\) .
\item
  One easily verifies that, for \(x\) , \(y\) in E and \(u \in A(\Phi)\) , one has \(\Phi(u(x), y) =: \Phi(x, \overline{u}(y))\) . Thus the adjoint endomorphism of a direct similarity \(u\) is the direct similarity \(\overline{u}\) conjugate of \(u\) in A( \(\Phi\) ).
\item
  Since every inverse similarity of E is the product of a direct similarity and the reflection with respect to \(Ae_{1}\) , the matrices of inverse similarities with respect to \((e_{1}, e_{2})\) are the matrices of the form \(\begin{pmatrix} a & -\delta b \\ b & -a \end{pmatrix}\) .
\end{enumerate}

We shall henceforth denote by S the group of similarities of E, by \(S^{+}\) the group of direct similarities, by H that of homotheties ≠ 0, and by \(O^{+}\) that of rotations. Recall that we have \(H \subset S^{+}\) (§ 6, no. 5).

Corollary 1. --- The group S \(^{+}\) of direct similarities is commutative. For any non-isotropic vectors x, y in E, there exists a unique direct similarity u such that y = u(x).

The first assertion follows from the fact that the algebra A(Φ) is commutative. Since E is a free monogenic A(Φ)-module generated by x (resp. y), there exists a unique element u (resp. u′) of A(Φ) such that \(y = u(x)\) (resp. \(x = u'(y)\) ); whence \(x = u(u'(x))\) , and \(uu'\) is the identity; this shows that u is invertible, and is therefore a direct similarity.

COROLLARY 2. --- The group O \(^{+}\) of rotations is commutative. For any vectors x, y in E such that Φ(x, x) = Φ(y, y) ≠ 0, there exists a unique rotation u such that y = u(x).

The first assertion follows from Cor. 1. This also shows that there exists a direct similarity \(u\) and exactly one such that \(y = u(x)\) ; since \(\Phi(u(x), u(x)) = \Phi(x, x)\) the multiplier of \(u\) is equal to 1, and \(u\) is therefore a rotation.

COROLLARY 3. --- The group S \(^{+}\) /H is commutative. It acts on the set of non-isotropic lines of E. Whatever the non-isotropic lines D, D′ of E may be, there exists one and only one element φ of S \(^{+}\) /H such that D' = φ(D).

This follows from Cor. 1 and from the fact that H leaves every line of E globally invariant.

PROPOSITION 2. --- The kernel of the canonical homomorphism from O \(^{+}\) into S \(^{+}\) /H is \{1, -1\}.

Indeed, 1 and -1 are the only elements of \(\mathbf{H} \cap \mathbf{O}^{+}\) .

PROPOSITION 3. --- The homomorphism \(u \to u/\overline{u} = u^{2}/\mathrm{N}(u)\) of \(\mathrm{S}^{+}\) into itself has H as its kernel, and defines, by passing to the quotient, an isomorphism from \(\mathrm{S}^{+}/\mathrm{H}\) onto \(\mathrm{O}^{+}\) .

Indeed, the relation \(u/\overline{u} = 1\) is equivalent to \(u = \overline{u}\) that is, to \(u \in A^* = H\) Thus H is the kernel of \(u \to u/\overline{u}\) . Since \(N(u/\overline{u}) = 1\) , \(u/\overline{u}\) is a rotation (Remark 1). It remains to show that every element \(v\) of \(O^+\) is of the form \(u/\overline{u}\) ( \(u \in S^+\) ). If \(1 + v\) is invertible, one can take \(u = 1 + v\) , because the relation \(N(v) = v\bar{v} = 1\) implies \(1 + v = v(1 + \bar{v})\) . Otherwise, one has \(N(1 + v) = (1 + v)(1 + \bar{v}) = 0\) that is, setting \(v = a + bw\) ( \(a \in A\) , \(b \in A\) , \(w^2 = \delta \in A\) ), \(2(1 + a) = 0\) , whence \(a = -1\) ; but the relations \(a = -1\) and \(N(v) = a^2 - \delta b^2 = 1\) imply \(b = 0\) , whence \(v = -1\) ; since \(\bar{w} = -w\) , it suffices, in this case, to take \(u = w\) .

When A(Φ) is a field, prop. 3 is a special case of Hilbert's theorem (chap.~V, § 11, no. 5, th. 3).

COROLLARY. — Let \(i: \mathrm{O}^{+} \to \mathrm{S}^{+}/\mathrm{H}\) and \(d: \mathrm{S}^{+}/\mathrm{H} \to \mathrm{O}^{+}\) denote the homomorphisms defined in props. 2 and 3, and write the abelian groups \(\mathrm{O}^{+}\) and \(\mathrm{S}^{+}/\mathrm{H}\) additively; we have \(d(i(\theta)) = 2\theta\) for \(\theta \in \mathrm{O}^{+}\) and \(i(d(\varphi)) = 2\varphi\) for \(\varphi \in \mathrm{S}^{+}/\mathrm{H}\) .

Indeed, if \(\bar{\nu}\) is a rotation, we have \(\bar{\nu} = \nu^{-1}\) , whence \(d(i(\nu)) = \nu/\bar{\nu} = \nu^{2}\) . On the other hand, if \(\varphi \in S^{+}/H\) , \(\varphi\) is the class mod. H of a direct similarity \(u\) , and \(d(\varphi) = u/\bar{u} = u^{2}/N(u)\) is congruent to \(u^{2}\) mod. H; whence the second formula.

\subsection{Plane trigonometry.}

In this section, we choose a generator w of the algebra A(Φ) such that \(w^{2} \in A\) . Such a generator is determined up to a homothety (no. 1, Remark 2), hence the element \(w^{2}\) of A, which we denote by δ, is determined modulo the multiplicative subgroup (A*)² of squares of nonzero elements of A.

Remark. — When -1 belongs to the class mod. (A*)² in question, one generally chooses w such that \(w^{2} = -1\) , which determines it up to sign. When this class contains 1 but not -1, one generally chooses w such that \(w^{2} = 1\) , which again determines it up to sign.

This being so, every element \(\nu\) of \(S^{+}\) can be written, in one and only one way, in the form

\[
\nu = c _ {w} (\nu) + s _ {w} (\nu) \varpi\tag{1}
\]

where \(c_w(v)\) , \(s_w(v)\) belong to A; the element \(s_w(v)/c_w(v)\) of the projective field \(\tilde{\mathbf{A}}\) (chap.~II, 2 \(^{e}\) ed., App. III, no. 5) is denoted \(t_w(v)\) ; it depends only on the class of \(v\) mod. H; thus \(t_w\) defines, by passage to the quotient, a map from S \(^{+}\) /H into the projective field \(\tilde{\mathbf{A}}\) , which we shall still denote by \(t_w\) by abuse of language. We shall often write \(c\) , \(s\) and \(t\) instead of \(c_w\) , \(s_w\) and \(t_w\) . We have \(c_{-w} = c_w\) , \(s_{-w} = -s_w\) and \(t_{-w} = -t_w\) .

PROPOSITION 4. — a) When \(w^2 = \delta\) is not a square in A (that is, when E contains no isotropic line), the map t from S⁺/H into \(\tilde{\mathbf{A}}\) is a bijection.

\begin{enumerate}
\def\labelenumi{\alph{enumi})}
\setcounter{enumi}{1}
\item
  When δ is the square of an element γ of A, the map t is a bijection from S⁺/H onto \(\tilde{\mathbf{A}}\) minus the elements 1/γ and -1/γ.
\item
  Writing \(\mathrm{S}^{+}/\mathrm{H}\) additively, we have , for \(\varphi, \varphi'\) in \(\mathrm{S}^{+}/\mathrm{H}\)
\end{enumerate}

\[
t (\varphi + \varphi^ {\prime}) = (t (\varphi) + t (\varphi^ {\prime})) / (1 + \delta t (\varphi) t (\varphi^ {\prime}))\tag{2}
\]

when \(t(\varphi)\) and \(t(\varphi')\) are finite and \(1 + t(\varphi)t(\varphi')\) is \(\neq 0\) .

Indeed, as \(S^{+}/H\) is a set of lines (with 0 removed) of \(A(\Phi)\) considered as a vector plane over A, t is injective. On the other hand, for an element \(a + bw\) ( \(a \in A\) , \(b \in A\) ) of \(A(\Phi)\) to be a direct similarity, it is necessary and sufficient that it be invertible, that is, that one have \(N(a + bw) = a^{2} - \delta b^{2} \neq 0\) , or again \((b/a)^{2} \neq 1/\delta\) ; this proves the surjectivity assertions in a) and b). Finally, the product of the similarities \(1 + t(\varphi)w\) and \(1 + t(\varphi')w\) is the similarity \(1 + \delta t(\varphi)t(\varphi') + (t(\varphi) + t(\varphi'))w\) , which proves c).

PROPOSITION 5. --- Let us write O \(^{+}\) additively. For every pair of elements θ, θ' of O \(^{+}\) we have

\begin{enumerate}
\def\labelenumi{(\arabic{enumi})}
\setcounter{enumi}{2}
\tightlist
\item
\end{enumerate}

\[
c (\theta) ^ {2} - \delta s (\theta) ^ {2} = 1\tag{4}
\]

\[
c (\theta + \theta^ {\prime}) = c (\theta) c (\theta^ {\prime}) + \delta s (\theta) s (\theta^ {\prime})\tag{5}
\]

\[
s (\theta + \theta^ {\prime}) = s (\theta) c (\theta^ {\prime}) + c (\theta) s (\theta^ {\prime}).
\]

The relation (3) indeed expresses that \(\mathrm{N}(c(\theta)+s(\theta)\omega)=1\) To prove (4) and (5), it suffices to compute, in \(\mathrm{A}(\Phi)\) the product of the rotations \(c(\theta)+s(\theta)\omega\) and \(c(\theta')+s(\theta')\omega\) .

PROPOSITION 6. — Let d be the isomorphism from S+/H onto O+ defined in Prop. 3. For every element φ of S+/H such that t = t(φ) is finite, one has

\[
s (d (\varphi)) = 2 t / (1 - \delta t ^ {2}), \quad c (d (\varphi)) = (1 + \delta t ^ {2}) / (1 - \delta t ^ {2}).\tag{6}
\]

Indeed \(\varphi\) is the class mod. H of the similarity \(1 + tw\) , and the rotation \(d(\varphi)\) is therefore \((1 + tw)^2/\mathrm{N}(1 + tw) = (1 + \delta t^2 + 2tw)/(1 - \delta t^2)\) (prop. 3), which proves (6).

COROLLARY. --- For every element θ of O⁺ such that \(t(\theta)\) is finite, we have

\[
s (2 \theta) = 2 t (\theta) / (1 - \delta t (\theta) ^ {2}), \quad c (2 \theta) = (1 + \delta t (\theta) ^ {2}) / (1 - \delta t (\theta) ^ {2}). \tag {7}
\]

For every element \(\varphi\) of \(S^{+}/H\) such that \(t(\varphi)\) is finite and \(1 + \delta t(\varphi)^{2} \neq 0\) , one has

\[
t (2 \varphi) = 2 t (\varphi) / (1 + \delta t (\varphi) ^ {2}).\tag{8}
\]

Indeed this follows immediately from prop. 6 and the corollary of prop. 3.

Remark. --- Formulas (6) remain true for \(t = \infty\) provided that the rational functions appearing on the right-hand side are replaced by their canonical extensions to the projective field. \(\tilde{\mathbf{A}}\) (chap.~II, 2 \(^{\text{e}}\) ed., App. III, no. 5); indeed, if \(t = \infty\) , \(\varphi\) is the class of \(w\) , and we have \(d(\varphi) = -1\) , \(s(d(\varphi)) = 0\) and \(c(d(\varphi)) = -1\) ; these are indeed the values taken by the canonical extensions of the right-hand sides for \(t = \infty\) The same holds for (7) when \(t(\theta) = \infty\) , and for (8) when \(t(\varphi) = \infty\) or that \(1 + \delta t(\varphi)^{2} = 0\) . Similarly, formula (2) remains true when only one of the elements \(t(\varphi)\) , \(t(\varphi')\) , \(t(\varphi)\) for example, is infinite, provided that its right-hand side is regarded as a rational function of \(t(\varphi)\) alone: indeed, the product of the similarities \(1 + t(\varphi')w\) and \(w\) is \(\delta t(\varphi') + w\) , while the value taken by the canonical extension of the right-hand side of (2) is \(1/\delta t(\varphi')\) . Finally, when \(t(\varphi)\) and \(t(\varphi')\) are finite and one has \(1 + \delta t(\varphi)t(\varphi') = 0\) , one has \(t(\varphi) + t(\varphi') = 0\) (otherwise \(t(\varphi)^{2}\) would be equal to \(1/\delta\) , which is impossible (prop. 4)); one may therefore agree that the value of the right-hand side of (2) is \(\infty\) , and this value is indeed that of the left-hand side. When \(t(\varphi)\) and \(t(\varphi')\) are both infinite, the right-hand side of (2) is not defined.

\subsection{Angles.}

We shall assume, in this section and the next, that A is an ordered field, hence of characteristic zero. Recall (Rectifications to Chap.~VI) that, if F is a vector space over A, the relation “ there exists λ \textgreater{} 0 such that y = λx ” is an equivalence relation between elements x, y of F - \{0\}, that every equivalence class for this relation is called an open half-line with origin 0, and that the union of an open half-line and \{0\} is called a closed half-line (or simply half-line) with origin 0; if D is a line and Δ a closed half-line contained in D, D is the union of Δ and -Δ, and contains no other closed half-line. We shall say that a half-line is isotropic if the line containing it is isotropic.

Recall also (ibid.) that, given a vector space of finite dimension n over A, an orientation on F is the choice of one of the two half-lines of the space \(\bigwedge^{n}F\) the n-vectors belonging to this half-line are said to be positive. A finite-dimensional vector space equipped with an orientation is said to be oriented.

The homotheties of E whose ratio is \textgreater0 obviously form a subgroup of index 2 in H; we will denote it by H \(^{+}\) . The canonical homomorphism i : O \(^{+}\) → S \(^{+}\) /H (cf.~prop. 2) is the composite of the canonical homomorphisms of S \(^{+}\) /H \(^{+}\) on S \(^{+}\) /H and of O \(^{+}\) in S \(^{+}\) /H \(^{+}\) ; since O \(^{+}\) ∩ H \(^{+}\) = \{1\} this latter homomorphism is injective.

PROPOSITION 7. --- Suppose that A is a maximal ordered field and that \(\Phi\) be a positive form (§ 7). Then the canonical homomorphisms of O \(^{+}\) into S \(^{+}\) /H \(^{+}\) and of O \(^{+}\) /\{1, -1\} into S \(^{+}\) /H are bijective, and S \(^{+}\) is isomorphic to O \(^{+}\) × H \(^{+}\) .

We have already seen that the homomorphisms in question are injective, and it suffices to show that the first one is surjective. Let \((e_{1}, e_{2})\) an orthonormal basis of E, and let w be the direct similarity such that \(w(e_{1}) = e_{2}\) (cor. 1 of prop. 1, no. 1); we then have \(w^{2} = -1\) (prop. 1, b)). Given any direct similarity transformation \(u = a + bw\) ( \(a \in \mathbf{A}, b \in \mathbf{A}\) ), we have \(\mathrm{N}(u) = a^2 + b^2 > 0\) and there exists one and only one rotation contained in the same half-line of \(\mathbf{A}(\Phi)\) as \(u\) , namely \((a^2 + b^2)^{1/2} u\) . QED.

COROLLARY. --- Given two half-lines D, D' with origin 0, there exists one and only one rotation v such that v(D) = D'.

The hypotheses indeed imply that E contains no isotropic lines. Our assertion then follows from cor. 1 of prop. 1 (no. 1).

We shall henceforth assume that A is a maximal ordered field, and that the form \(\Phi\) is positive. In the set of pairs \((\mathrm{D}_1, \mathrm{D}_2)\) of lines (resp. half-lines with origin 0) of E, the relation “there exists a direct similarity (resp. a rotation) u such that \(u(\mathrm{D}_1) = \mathrm{D}_1'\) and \(u(\mathrm{D}_2) = \mathrm{D}_2'\) ” is an equivalence relation between pairs \((\mathrm{D}_1, \mathrm{D}_2)\) and \((\mathrm{D}_1', \mathrm{D}_2')\) The equivalence class of the pair \((\mathrm{D}_1, \mathrm{D}_2)\) is called, by definition, the angle of the lines (resp. half-lines) \(\mathrm{D}_1, \mathrm{D}_2\) (taken in this order); we denote it \((\widehat{\mathrm{D}_1, \mathrm{D}_2})\) .

PROPOSITION 8. --- Suppose that A is a maximal ordered field, and that the form \(\Phi\) is positive. Let \(D_{1}, D_{2}, D_{1}^{\prime}, D_{2}^{\prime}\) be four lines (resp. half-lines) with origin 0 in E. For the angles \((\widehat{D_{1}, D_{2}})\) and \((\widehat{D_{1}^{\prime}, D_{2}^{\prime}})\) to be equal, it is necessary and sufficient that the angles \((\widehat{D_{1}, D_{1}^{\prime}})\) and \((\widehat{D_{2}, D_{2}^{\prime}})\) be equal.

Let us prove the necessity of the stated condition. Let \(u\) be a direct similarity (resp. a rotation) such that \(u(\mathrm{D}_{1}) = \mathrm{D}_{1}^{\prime}\) and \(u(\mathrm{D}_{2}) = \mathrm{D}_{2}^{\prime}\) . By cor. 3 of prop. 1, there exists a direct similarity (resp. by the corollary of prop. 7, a rotation) \(\nu\) such that \(\nu(\mathrm{D}_{1}) = \mathrm{D}_{2}\) . Since the group \(\mathrm{S}^{+}\) (resp. \(\mathrm{O}^{+}\) ) is commutative, we have \(\mathrm{D}_{2}^{\prime} = u(\nu(\mathrm{D}_{1})) = \nu(u(\mathrm{D}_{1})) = \nu(\mathrm{D}_{1}^{\prime})\) , and this shows that \((\widehat{\mathrm{D}_{1}, \mathrm{D}_{1}^{\prime})} = (\widehat{\mathrm{D}_{2}, \mathrm{D}_{2}^{\prime}})\) . Sufficiency follows from necessity by interchanging \(\mathrm{D}_{2}\) and \(\mathrm{D}_{1}^{\prime}\) .

It follows from prop. 8 that, to every angle \((\widehat{\mathrm{D}_1,\mathrm{D}_2})\) of lines (resp. half-lines) with origin 0 in E, there is canonically associated a well-determined element of \(\mathrm{S}^{+}/\mathrm{H}\) (resp. \(\mathrm{O}^{+}\) ), namely the class mod. H of the direct similarities \(\nu\) (resp. the rotation \(\nu\) ) such that \(u(D_1) = D_2\) for any representative ( \(D_1, D_2\) ) of the angle ( \(\widehat{D_1, D_2}\) ). We have thus defined a canonical bijection \(h\) (resp. \(h'\) ) from the set \(\mathfrak{A}_0\) of angles of lines (resp. \(\mathfrak{A}\) of angles of half-lines) onto \(S^+/H\) (resp. \(O^+\) ); in particular, for every \(\varphi \in \mathfrak{A}\) , one says that \(h(\varphi)\) is the rotation with angle \(\varphi\) . We shall transport to \(\mathfrak{A}_0\) and \(\mathfrak{A}\) , by means of \(h^{-1}\) and of \(h'^{-1}\) , the structures of commutative groups of \(S^+/H\) and of \(O^+\) , and we shall denote additively the groups \(\mathfrak{A}_0\) and \(\mathfrak{A}\) thus obtained. If one denotes by D, \(D', D''\) lines (resp. half-lines) with origin 0 in E, one has by definition

\[
\widehat {(\mathbf {D} , \mathbf {D} ^ {\prime \prime})} = \widehat {(\mathbf {D} , \mathbf {D} ^ {\prime})} + \widehat {(\mathbf {D} ^ {\prime} , \mathbf {D} ^ {\prime \prime})}\tag{9}
\]

(Chasles relation);

one deduces from this

\[
\widehat {(\mathrm{D} , \mathrm{D})} = 0, \quad \widehat {(\mathrm{D} , \mathrm{D} ^ {\prime})} = - (\widehat {\mathrm{D} ^ {\prime} , \mathrm{D}}).\tag{10}
\]

Remarks. --- 1) The set L of lines (resp. half-lines) with origin 0 in E is a homogeneous space of the abelian group S+/H (resp. O+) such that the identity element is the only operator leaving all elements of L invariant. One can therefore apply to L the formulas of chap.~II, 2 \(^{e}\) ed., App. II, n° 1; prop. 8 is thus a particular case of the "parallelogram rule", and formulas (9) and (10) are particular cases of formulas (2) (ibid.).

\begin{enumerate}
\def\labelenumi{\arabic{enumi})}
\setcounter{enumi}{1}
\tightlist
\item
  In the definition of the group of angles of lines, one can, instead of the group S+/H, use the group O+/\{−1, 1\} which is canonically isomorphic to it (prop. 7). The canonical homomorphism of O+ onto O+/\{−1, 1\} thus corresponds to a homomorphism of A onto A0, namely the one which, to the angle of the two half-lines Δ, Δ', associates the angle of the lines D, D' containing respectively Δ, Δ'. *In the case where the field A is the field of real numbers, the group A is thus a double covering of the group A0.*
\end{enumerate}

By prop. 8, all angles of lines (resp. half-lines) of the form \((\widehat{\mathbf{D}^{\prime},\mathbf{D}^{\prime\prime}})\) where \(\mathbf{D}'\) and \(\mathbf{D}''\) are orthogonal (resp.

of the form \((\widehat{\mathrm{D}}, -\mathrm{D})\) are equal: this indeed follows from Remark 2 of No. 1 (resp. is obvious). This angle of lines (resp. of half-lines) is called the right angle (resp. the straight angle); it is an element of order 2 of \(A_{0}\) (resp. of A).

PROPOSITION 9. --- Suppose that the field A is maximally ordered, and that \(\Phi\) is a positive form. For every integer \(n>1\) , the number of elements \(\theta\) of the group \(\mathfrak{A}_{0}\) of angles of lines (resp. \(\mathfrak{A}\) of angles of half-lines) such that \(n\theta=0\) is equal to \(n\) .

As \(\mathfrak{A}_{0}\) , \(\mathrm{S}^{+}/\mathrm{H}\) , \(\mathfrak{A}\) and \(\mathrm{O}^{+}\) are isomorphic (prop. 3), it suffices to carry out the proof for \(\mathrm{O}^{+}\) , that is, to show that there are exactly \(n\) rotations \(\varrho\) such that \(\varphi^{n}=1\) Since A is a maximal ordered field and that \(\mathrm{A}(\Phi)\) is a degree 2 extension field of A (prop. 1 a)), \(\mathrm{A}(\Phi)\) is an algebraically closed field (chap.~VI, § 2, no. 6, th. 3). Hence the roots \(n\) -th roots of unity in \(\mathrm{A}(\Phi)\) form a cyclic group of order \(n\) (chap. V, § 11, no. 1, th. 1). Since we have \(\mathrm{N}(u)=u\overline{u}\geqslant0\) for all \(u\in\mathrm{A}(\Phi)\) , the relation \(u^{n}=1\) implies that one has \(\mathrm{N}(u)=1\) , hence that \(u\) is a rotation (No. 1, Remark 1). This proves our assertion.

COROLLARY. --- The right angle (resp. straight angle) is the only element of order 2 in the group. \(\mathfrak{A}_{0}\) (resp. \(\mathfrak{A}\) ).

We will assume finally that the plane E is oriented.

Lemma 1. --- Let u be a direct similarity of E; all bivectors of the form \(x \wedge u(x)\) belong to the same closed half-line of \(\bigwedge^{2}\mathbf{E}\) .

The case where \(u\) is a homothety is trivial. Otherwise, we have \(x \wedge u(x) \neq 0\) for all \(x \neq 0\) ; let \(x, y\) be two vectors of \(E\) ( \(x \neq 0, y \neq 0\) ); there exists \(v \in S^+\) such that \(y = v(x)\) , whence \(y \wedge u(y) = v(x) \wedge uv(x) = v(x) \wedge v(u(x)) = (\det v)(x \wedge u(x))\) ; by taking an orthonormal basis of \(E\) , we see that det \(v\) is positive (prop. 1 b)); hence our assertion.

That being said, among the two generators \(w\) of \(\mathrm{A}(\Phi)\) such that \(w^{2} = -1\) there exists one and only one such that the bivector \(x \wedge w(x)\) be positive for all \(x\in\mathbf{E}\) . It is this generator that we will choose to define the functions \(c_{w}, s_{w}\) and \(t_{w}\) (no. 2). Let \(h\) and \(h'\) the canonical bijections defined above from the group \(\mathfrak{A}_{0}\) angles of lines on S+/H and of the group \(\mathfrak{A}\) of half-line angles on O+. The composed mappings \(t_{w}\circ h\) of \(\mathfrak{A}_{0}\) in the projective field \(\tilde{\mathbf{A}}\) , \(c_{w}\circ h'\) and \(s_{w}\circ h'\) of \(\mathfrak{A}\) in the field A are denoted respectively by tg, cos, and sin, and are called the tangent, cosine, and sine functions. The mapping \(\varphi\to1/\mathrm{tg}\varphi\) of \(\mathfrak{A}_{0}\) in \(\tilde{\mathbf{A}}\) is denoted cotg and is called the cotangent function. The functions cosine, sine, tangent, and cotangent are said to be the trigonometric functions. The composite mappings tg∘p and cotg∘p, where p denotes the canonical homomorphism from \(\mathfrak{A}\) on \(\mathfrak{A}_{0}\) (Remark 2 above) are still denoted tg and cotg by abuse of language.

Formulas (2), (8), (3), (4), (5), and (7) of §2 give, since here we have \(\delta = -1\)

\begin{enumerate}
\def\labelenumi{(\arabic{enumi})}
\setcounter{enumi}{10}
\tightlist
\item
\end{enumerate}

\[
\operatorname{tg} \left(\varphi + \varphi^ {\prime}\right) = (\operatorname{tg} \varphi + \operatorname{tg} \varphi^ {\prime}) / (1 - \operatorname{tg} \varphi \operatorname{tg} \varphi^ {\prime})\tag{12}
\]

\[
\operatorname{tg} (2 \varphi) = 2 \operatorname{tg} \varphi / (1 - \operatorname{tg} ^ {2} \varphi)
\]

for \(\varphi, \varphi'\) in \(\mathfrak{A}_0\) ;

\begin{enumerate}
\def\labelenumi{(\arabic{enumi})}
\setcounter{enumi}{12}
\tightlist
\item
\end{enumerate}

\[
\cos^ {2} \theta + \sin^ {2} \theta = 1\tag{14}
\]

\[
\cos (\theta + \theta^ {\prime}) = \cos \theta \cos \theta^ {\prime} - \sin \theta \sin \theta^ {\prime}\tag{15}
\]

\[
\sin (\theta + \theta^ {\prime}) = \sin \theta \cos \theta^ {\prime} + \cos \theta \sin \theta^ {\prime}\tag{16}
\]

\[
\sin (2 \theta) = 2 \operatorname{tg} \theta / (1 + \operatorname{tg} ^ {2} \theta),
\]

\[
\cos (2 \theta) = (1 - \mathrm{tg} ^ {2} \theta) / (1 + \mathrm{tg} ^ {2} \theta)
\]

for \(\theta, \theta'\) in \(\mathfrak{A}\) On the other hand, we have, by definition or as an easy consequence of the preceding formulas:

\begin{enumerate}
\def\labelenumi{(\arabic{enumi})}
\setcounter{enumi}{16}
\tightlist
\item
\end{enumerate}

\[
\operatorname{tg} \theta = \sin \theta / \cos \theta , \quad \cot \mathrm{g} \theta = \cos \theta / \sin \theta\tag{18}
\]

\[
1 + \mathrm{tg} ^ {2} \theta = 1 / \cos^ {2} \theta , \quad 1 + \cot \mathrm{g} ^ {2} \theta = 1 / \sin^ {2} \theta
\]

for θ ∈ 𝒜.

Given two nonzero vectors \(x, y\) of E, the angle between these two vectors (taken in this order) is called, and denoted \((\widehat{x}, y)\) the angle of the half-lines to which they belong. For every vector \(x\) from E we call the length of \(x\) is called, and denoted \(|x|\) , the element \(\Phi(x, x)^{1/2}\) of A.

PROPOSITION 10. --- Suppose that the field A is maximally ordered, that the plane E is oriented, and that the form \(\Phi\) is positive. For every pair of nonzero vectors \(x\) , \(y\) of E on \(a\)

\begin{enumerate}
\def\labelenumi{(\arabic{enumi})}
\setcounter{enumi}{18}
\tightlist
\item
\end{enumerate}

\[
\cos \widehat {(x , y)} = \Phi (x, y) / | x |. | y |\tag{20}
\]

\[
\widehat {\sin (x , y)}. e = (x \wedge y) / | x |. | y |,
\]

where e denotes the positive bivector such that \(\Phi_{(2)}(e, e) = 1\) .

Indeed, since the vectors \(x' = x / |x|\) and \(y' = y / |y|\) are both of length 1, there exists a rotation \(\nu\) and exactly one such that \(\nu(x') = y'\) (No. 1, cor. 2 of prop. 1). If we set \(\nu = a + bw\) ( \(a, b\) In A), we have by definition \(a = \cos(\widehat{x,y})\) and \(b = \sin(\widehat{x,y})\) . The relation \(y' = \nu(x') = ax' + bw(x')\) gives \(\Phi(x', y') = a\Phi(x', x') = a\) since \(x'\) and \(w(x')\) are orthogonal (Remark 2 of No. 1); this proves (19). On the other hand, this relation also gives \(x' \wedge y' = bx' \wedge w(x') = b.e\) by the definition of the extension \(\Phi_{(2)}\) of \(\Phi\) à \(\bigwedge^{2} E\) (§ 1, no. 9, formula (37)) and the choice of \(w\) ; this proves (20).

Remarks. --- 3) Given two lines (resp. rays) D, D' of an affine plane L attached to E, one calls the angle of D and D', denoted (D, D'), the angle made by their directions in E (resp. the rays with origin 0 of E corresponding to D and D') (chap.~II, 2). \(^{e}\) ed., App. II, nos. 1 and 3).

\begin{enumerate}
\def\labelenumi{\arabic{enumi})}
\setcounter{enumi}{3}
\tightlist
\item
  Let F be a vector space of arbitrary dimension over the maximal ordered field A, and \(\Psi\) a positive non-degenerate symmetric bilinear form on F. Given two vectors \(x\) , \(y\) linearly independent elements of F, let \(\mathbf{F}'\) the vector plane they span; we call the angle of \(x\) and \(y\) the angle of \(x\) and \(y\) considered as elements of the plane \(\mathbf{F}'\) ; it is denoted \((x,y)\) The cosine of this angle is, by virtue of (19), given by
\end{enumerate}

\[
\cos \widehat {(x , y)} = \Psi (x, y) / | x |. | y |\tag{21}
\]

\((\text{ou } |x| = \Psi(x, x)^{1/2}\) is also called the length of the vector \(x)\) and is therefore independent of the chosen orientation on \(F'\) ; the sine and the tangent of \((x, y)\) change sign if one changes the orientation of \(F'\) Given two nonzero proportional vectors \(x, y\) of \(F\) , set \((x, y) = 0\) by convention.
